\documentclass[11pt,a4paper]{article}
\usepackage[utf8]{inputenc}
\usepackage[T1]{fontenc}
\usepackage[french]{babel}
\usepackage{amsmath,amssymb,mathtools}
\usepackage{geometry}
\usepackage{graphicx}
\usepackage{booktabs}
\usepackage{hyperref}
\geometry{margin=2.4cm}
\hypersetup{colorlinks=true,linkcolor=blue!50!black,urlcolor=blue!50!black}

\newcommand{\dd}{\mathrm{d}}
\newcommand{\lp}{\ell_{\mathrm P}}
\newcommand{\heff}{h_{\mathrm{eff}}}
\newcommand{\J}{\mathcal{J}}
\newcommand{\Ocal}{\mathcal{O}}

\title{Implications candidates d'une cellule de phase deformee\\\large Version prepublication de travail}
\author{Document de travail}
\date{22 mai 2026}

\begin{document}
\maketitle

\begin{abstract}
On etudie une mesure de phase deformee
\[
\dd N=\frac{g\,\dd^3x\,\dd^3p}{h^3(1+\beta p^2)^\gamma}
\]
et en particulier le cas isotrope $\gamma=3$. Le modele donne une densite finie de modes, une regularisation UV de l'energie du vide, des bornes pour les secteurs a occupation bornee, et des corrections au rayonnement noir. Une conclusion importante est que les bosons thermiques ne possedent pas d'energie maximale : a haute temperature, le rayonnement passe de $u\propto T^4$ a $u\propto T$.
\end{abstract}

\section{Hypothese de mesure}
La cellule effective est
\[
\heff^3(p)=h^3(1+\beta p^2)^\gamma.
\]
Le cas central est
\[
\J(p)=(1+\beta p^2)^3.
\]
Cette mesure doit etre comprise comme un ansatz de modele, ou comme la limite phenomenologique d'une algebre deformee precise. Le choix $p^2=p_x^2+p_y^2+p_z^2$ fixe un repere et n'est pas covariant par lui-meme.

\section{Integrale maitresse}
Pour
\[
I_s(\gamma)=\int_0^\infty \frac{p^{s-1}\,\dd p}{(1+\beta p^2)^\gamma},
\]
le changement de variable $u=\beta p^2$ donne
\[
I_s(\gamma)=\frac12\beta^{-s/2}B\left(\frac{s}{2},\gamma-\frac{s}{2}\right),\qquad \gamma>\frac{s}{2}.
\]
Donc, pour $\gamma=3$,
\[
\int_0^\infty \frac{p^2\,\dd p}{(1+\beta p^2)^3}=\frac{\pi}{16\beta^{3/2}},\qquad
\int_0^\infty \frac{p^3\,\dd p}{(1+\beta p^2)^3}=\frac{1}{4\beta^2}.
\]

\section{Densite finie de modes}
La densite maximale de modes monoparticulaires est
\[
 n_{\max}=\frac{4\pi g}{h^3}\int_0^\infty \frac{p^2\,\dd p}{(1+\beta p^2)^3}
=\frac{g\pi^2}{4h^3\beta^{3/2}}.
\]
Avec $\beta=\beta_0\lp^2/\hbar^2$,
\[
 n_{\max}=\frac{g}{32\pi\beta_0^{3/2}\lp^3}.
\]
Le cumul exact, avec $q=\sqrt{\beta}p$, est
\[
N(p)=\frac{g\pi V}{2h^3\beta^{3/2}}\left[\arctan q-\frac{q(1-q^2)}{(1+q^2)^2}\right].
\]

\section{Rayonnement noir}
Pour $E=pc$,
\[
u(E,T)=\frac{4\pi g}{h^3c^3}\frac{E^3}{(e^{E/k_BT}-1)(1+\beta E^2/c^2)^\gamma}.
\]
Le maximum spectral verifie, avec $x=E/k_BT$ et $\epsilon=\beta(k_BT)^2/c^2$,
\[
3-\frac{x e^x}{e^x-1}-\frac{2\gamma\epsilon x^2}{1+\epsilon x^2}=0.
\]
A premier ordre,
\[
x_{\max}=2.82143937-18.2282675\,\gamma\,\epsilon+\Ocal(\epsilon^2).
\]
La correction de Stefan-Boltzmann est
\[
u(T)=aT^4\left[1-\gamma\frac{40\pi^2}{21}\frac{\beta k_B^2T^2}{c^2}+\Ocal(\beta^2T^4)\right],\qquad
 a=\frac{4\pi^5gk_B^4}{15h^3c^3}.
\]
En revanche, a tres haute temperature,
\[
u(T)\sim n_{\max}k_BT.
\]
C'est une correction conceptuelle majeure : le nombre de modes est fini, mais l'occupation bosonique n'est pas bornee.

\section{Fermions degeneres et vide}
Pour un gaz fermionique ultrarelativiste a temperature nulle,
\[
\rho(q_F)=\frac{g\pi c}{h^3\beta^2}\frac{q_F^4}{(1+q_F^2)^2},\qquad
\rho_{\max}=\frac{g\pi c}{h^3\beta^2}.
\]
Si $E=pc$, alors $P=\rho/3$ et
\[
P_{\max}=\frac{g\pi c}{3h^3\beta^2}.
\]
Pour un champ massless,
\[
\rho_{\mathrm{vac}}=\frac{4\pi g}{h^3}\int_0^\infty \frac{p^2\,\dd p}{(1+\beta p^2)^3}\frac12 pc
=\frac{g\pi c}{2h^3\beta^2}.
\]
Cette quantite est finie, mais ne resout pas seule le probleme de la constante cosmologique.

\section{Cosmologie effective}
Une dynamique de rebond peut etre postulee sous la forme
\[
H^2=\frac{8\pi G}{3}\rho\left(1-\frac{\rho}{\rho_{\max}}\right).
\]
Elle ne decoule pas automatiquement de la mesure de phase. Avec $w$ constant et conservation standard,
\[
a(t)=a_b\left[1+6\pi G\rho_{\max}(1+w)^2t^2\right]^{1/[3(1+w)]}.
\]
Une version publiable doit etudier les perturbations, la stabilite, la covariance et les contraintes observationnelles.

\section{Programme de publication}
Les points a developper sont : derivation de la mesure de Liouville depuis une algebre precise ; formulation covariante ; equation d'etat massive et interactive pour objets compacts ; calcul d'entropie avec contrainte gravitationnelle ; confrontation aux bornes sur $\beta_0$.

\end{document}
